% Part: turing-machines
% Chapter: machines-computations
% Section: representing-tms

\documentclass[../../../include/open-logic-section]{subfiles}

\begin{document}

\olfileid{tur}{mac}{rep}
\olsection{表示Turing机}

\begin{explain}
Turing机可以用\emph{状态图}来直观地表示。
状态图由状态节点和连接它们的箭头构成。自然地，每个状态节点代表机器的一个状态。
每个箭头代表从该状态可以执行的一条指令，指令的具体内容写在该箭头的上方或下方。
考虑下面这台机器，它只有两个内部状态 $q_0$ 和 $q_1$，以及一条指令：
\[
\begin{tikzpicture}[->,>=stealth',shorten >=1pt,auto,node distance=2.8cm,
                    semithick]
  \tikzstyle{every state}=[fill=none,draw=black,text=black]

  \node[initial,state]   (A)              {$q_0$};
  \node[state]   (B) [right of=A] {$q_1$};

  \path (A) edge  node {\TMtrans{\TMblank}{\TMstroke}{\TMright}} (B);
\end{tikzpicture}
\]
回想一下，Turing机有一个读写头和一个写有输入的纸带。
这条指令可以读作：\emph{如果在状态 $q_0$ 下扫描到~$\TMblank$，则写下一个~$\TMstroke$，向右移动，并转移到状态 $q_1$}。
这等价于将 $\tuple{q_0, \TMblank}$ 映射到 $\tuple{q_1, \TMstroke, \TMright}$ 的转移函数。
\end{explain}

\begin{ex}
\emph{偶数机}：下述Turing机停机，当且仅当纸带上有偶数个 $\TMstroke$（假设所有 $\TMstroke$ 都位于纸带上第一个 $\TMblank$ 之前）。
\[
\begin{tikzpicture}[->,>=stealth',shorten >=1pt,auto,node distance=2.8cm,
                    semithick]
  \tikzstyle{every state}=[fill=none,draw=black,text=black]

  \node[initial,state]         (A)              {$q_0$};
  \node[state]         (B) [right of=A] {$q_1$};

  \path (A) edge [bend left] node {\TMtrans{\TMstroke}{\TMstroke}{\TMright}} (B)
        (B) edge [loop above] node {\TMtrans{\TMblank}{\TMblank}{\TMright}} (B)
            edge [bend left] node {\TMtrans{\TMstroke}{\TMstroke}{\TMright}} (A);
\end{tikzpicture}
\]

该状态图对应于如下转移函数：
\begin{align*}
\delta(q_0, \TMstroke) & = \tuple{q_1, \TMstroke, \TMright},\\
\delta(q_1, \TMstroke) & = \tuple{q_0, \TMstroke, \TMright},\\
\delta(q_1, \TMblank)  & = \tuple{q_1, \TMblank, \TMright}
\end{align*}
\end{ex}

\begin{explain}
上述机器仅当输入是偶数个 $\TMstroke$ 时才停机。
否则，该机器（理论上）将无限地继续运行。
对于任何机器和输入，都可以追踪机器的\emph{格局}以确定其输出。
我们稍后会给出格局的形式定义。现在，我们可以直观地将格局看作一系列显示机器在运行中任意时刻状态的图示。
格局显示了纸带的内容、机器的状态以及读写头的位置。

让我们来追踪偶数机在输入为四个 $\TMstroke$ 时的格局。在这种情况下，我们预计机器会停机。
接着，我们再在输入为三个 $\TMstroke$ 的情况下运行该机器，此时机器将永远运行下去。

机器从状态~$q_0$ 开始，扫描最左边的~$\TMstroke$。
我们可以将机器的初始格局表示如下：
\[
\TMendtape \TMstroke_0 \TMstroke \TMstroke \TMstroke \TMblank \ldots
\]
上面的格局很直观。可以看到，机器从状态 $q_0$ 开始，扫描最左边的~$\TMstroke$。
这通过在最左边的~$\TMstroke$ 上标注状态名的下标来表示。
此时适用的指令是
$\delta(q_0, \TMstroke) = \tuple{q_1, \TMstroke, \TMright}$，因此
机器在纸带上向右移动并切换到状态~$q_1$。
\[
\TMendtape \TMstroke \TMstroke_1 \TMstroke \TMstroke \TMblank \ldots
\]
由于机器现在处于状态~$q_1$ 并扫描到一个~$\TMstroke$，我们必须“遵循”指令
$\delta(q_1, \TMstroke) = \tuple{q_0,
  \TMstroke, \TMright}$。这产生了格局
\[
\TMendtape \TMstroke \TMstroke \TMstroke_0 \TMstroke \TMblank \ldots
\]
随着机器的继续运行，规则按同样的顺序再次被应用，产生了以下两个格局：
\[
\TMendtape \TMstroke \TMstroke \TMstroke \TMstroke_1 \TMblank \ldots
\]
\[
\TMendtape \TMstroke \TMstroke \TMstroke \TMstroke \TMblank_0 \ldots
\]
机器现在处于状态~$q_0$，扫描到一个~$\TMblank$。
根据转移图，我们可以很容易地看到没有要执行的指令，因此机器已经停机。
这意味着输入被接受了。

接下来，假设我们以三个 $\TMstroke$ 作为输入启动机器。
最初的几个格局是类似的，因为执行的是相同的指令，只是纸带输入有微小差异：
\[
\TMendtape \TMstroke_0 \TMstroke \TMstroke \TMblank \ldots
\]
\[
\TMendtape \TMstroke \TMstroke_1 \TMstroke \TMblank \ldots
\]
\[
\TMendtape \TMstroke \TMstroke \TMstroke_0 \TMblank \ldots
\]
\[
\TMendtape \TMstroke \TMstroke \TMstroke \TMblank_1 \ldots
\]
机器现在已经扫描完所有的 $\TMstroke$，并在状态~$q_1$ 下读取到一个~$\TMblank$。
如图所示，存在一条形式为
$\delta(q_1, \TMblank) =\tuple{q_1,
\TMblank, \TMright}$ 的指令。
由于纸带向右无限延伸的部分都填满了 $\TMblank$，机器将\emph{永远}继续执行这条指令，
保持在状态~$q_1$ 并不断向右移动。
机器将永远不会停机，并且不接受该输入。
\end{explain}

\begin{explain}
需要注意的是，并非所有机器都会停机。
如果停机意味着机器没有指令可执行，那么我们只需确保每个状态下对每个符号都有一个发出的箭头，
就可以轻易构造出一台永不停机的机器。
通过为在状态~$q_0$ 下扫描到 $\TMblank$ 的情况添加一条指令，可以修改偶数机使其无限运行下去。
\end{explain}

\begin{ex}
\[
\begin{tikzpicture}[->,>=stealth',shorten >=1pt,auto,node distance=2.8cm,
                    semithick]
  \tikzstyle{every state}=[fill=none,draw=black,text=black]

  \node[initial,state] (A)              {$q_0$};
  \node[state]         (B) [right of=A] {$q_1$};

  \path
  (A) edge [bend left] node {\TMtrans{\TMstroke}{\TMstroke}{\TMright}} (B)
      edge [loop above] node {\TMtrans{\TMblank}{\TMblank}{\TMright}} (A)
  (B) edge [loop above] node {\TMtrans{\TMblank}{\TMblank}{\TMright}} (B)
      edge [bend left] node {\TMtrans{\TMstroke}{\TMstroke}{\TMright}} (A);
\end{tikzpicture}
\]
\end{ex}

\begin{explain}
机器表是表示Turing机的另一种方法。机器表将纸带字母表展示在 $x$ 轴上，而将机器状态的集合展示在 $y$ 轴上。在表格内，每个状态与符号交叉处的单元格内，写的是该指令的其余部分——新的状态、新的符号以及移动方向。机器表让人容易确定机器在何种状态下、扫描到何种符号时会停机。凡在表格中出现空缺的地方，就是机器可能的一个停机点。与状态图和指令集不同（在那里停机点并非总是一目了然），任何停机点都可以通过寻找机器表中的空缺位置而快速识别出来。
\end{explain}

\begin{ex}
偶数机的机器表是：
\[
\centering
\begin{tabular}{lllll}
\cline{2-4}
\multicolumn{1}{l|}{}      & \multicolumn{1}{c|}{$\TMblank$}
& \multicolumn{1}{c|}{$\TMstroke$}     & \multicolumn{1}{c|}{$\TMendtape$}   &  \\ \cline{1-4}
\multicolumn{1}{|l|}{$q_0$} & \multicolumn{1}{l|}{}
& \multicolumn{1}{l|}{\TMtrans{\TMstroke}{q_1}{\TMright}} &
\multicolumn{1}{l|}{\phantom{\TMtrans{\TMstroke}{q_1}{\TMright}}} &  \\ \cline{1-4}
\multicolumn{1}{|l|}{$q_1$} & \multicolumn{1}{l|}{\TMtrans{\TMblank}{q_1}{\TMright}}
& \multicolumn{1}{l|}{\TMtrans{\TMstroke}{q_0}{\TMright}} &
\multicolumn{1}{l|}{\phantom{\TMtrans{\TMstroke}{q_1}{\TMright}}} &  \\ \cline{1-4}
\end{tabular}
\]
如我们所见，当在状态~$q_0$ 下扫描到一个~$\TMblank$ 时，该机器停机。
\end{ex}

\begin{explain}
到目前为止，我们只考虑了能读取并接受输入的机器。然而，Turing机同时具备读取和写入的能力。这类机器的一个例子（尽管这样的例子非常多）是\emph{倍乘器}。一个倍乘器，当以纸带上一块包含~$n$ 个 $\TMstroke$ 的输入启动时，会输出一块包含~$2n$ 个 $\TMstroke$ 的符号。
\end{explain}

\begin{ex}
  \ollabel{ex:doubler} 在构建倍乘器之前，想出一个解决问题的\emph{策略}是很重要的。由于机器（按我们的形式化设定）无法记住它已读了多少个 $\TMstroke$，我们需要想出一种方法来跟踪纸带上所有的 $\TMstroke$。一种方法是用一个~$\TMblank$ 将输出区与输入区分隔开。然后，机器可以擦除输入中的第一个 $\TMstroke$，越过输入的其余部分，留下一个 $\TMblank$，并写下两个新的 $\TMstroke$。接着，机器将返回并找到输入中的第二个 $\TMstroke$，并同样将其加倍。对于输入中的每一个 $\TMstroke$，它将写出两个 $\TMstroke$ 作为输出。通过在机器运行时擦除输入，我们可以保证没有 $\TMstroke$ 被遗漏或重复加倍。当整个输入被擦除后，纸带上将剩下 $2n$ 个 $\TMstroke$。最终得到的Turing机的状态图如 \olref{fig:doubler} 所示。
  \begin{figure}
\[
\begin{tikzpicture}[->,>=stealth',shorten >=1pt,auto,node distance=2.8cm,
                    semithick]
  \tikzstyle{every state}=[fill=none,draw=black,text=black]

  \node[initial,state] (1)              {$q_0$};
  \node[state]         (2) [right of=1] {$q_1$};
  \node[state]         (3) [right of=2] {$q_2$};
  \node[state]         (4) [below of=3] {$q_3$};
  \node[state]         (5) [left of=4]  {$q_4$};
  \node[state]         (6) [left of=5]  {$q_5$};

  \path (1) edge node {\TMtrans{\TMstroke}{\TMblank}{\TMright}} (2)
    (2) edge [loop above] node {\TMtrans{\TMstroke}{\TMstroke}{\TMright}} (2)
      edge node {\TMtrans{\TMblank}{\TMblank}{\TMright}} (3)
    (3) edge [loop above] node {\TMtrans{\TMstroke}{\TMstroke}{\TMright}} (3)
        edge node {\TMtrans{\TMblank}{\TMstroke}{\TMright}} (4)
    (4) edge [loop below] node {\TMtrans{\TMblank}{\TMstroke}{\TMleft}} (4)
        edge node {\TMtrans{\TMstroke}{\TMstroke}{\TMleft}} (5)
    (5) edge [loop below]  node {\TMtrans{\TMstroke}{\TMstroke}{\TMleft}} (5)
        edge              node {\TMtrans{\TMblank}{\TMblank}{\TMleft}} (6)
    (6) edge [loop below] node {\TMtrans{\TMstroke}{\TMstroke}{\TMleft}} (6)
        edge              node {\TMtrans{\TMblank}{\TMblank}{\TMright}} (1);
\end{tikzpicture}
\]
\caption{一个倍乘器}
\ollabel{fig:doubler}
\end{figure}
\end{ex}

\begin{prob}
任选一个输入，并追踪 \olref[tur][mac][rep]{ex:doubler} 中倍乘器的格局。
\end{prob}

\begin{prob}
设计一个字母表为 $\{\TMendtape,\TMblank, A, B\}$ 的Turing机，使其接受（即在该输入上停机）任何由 $A$ 和 $B$ 组成的字符串，其中 $A$ 的数量等于 $B$ 的数量\emph{且}所有的 $A$ 都位于所有的 $B$ 之前；并且拒绝（即在该输入上不停机）任何 $A$ 的数量不等于 $B$ 的数量或者 $A$ 不全部位于 $B$ 之前的字符串。（例如，该机应接受 $AABB$ 和 $AAABBB$，但拒绝 $AAB$ 和 $AABBAABB$。）
\end{prob}

\begin{prob}
设计一个字母表为 $\{\TMendtape,\TMblank, A, B\}$ 的Turing机，它以由 $A$ 和 $B$ 组成的任意字符串 $\alpha$ 作为输入，并将其复制，从而产生一个形如 $\alpha\alpha$ 的输出。（例如，输入 $ABBA$ 应产生输出 $ABBAABBA$。）
\end{prob}

\begin{prob}
\emph{按字母顺序？：}设计一个字母表为 $\{\TMendtape,\TMblank, A, B\}$ 的Turing机，当给定一个由有限个 $A$ 和 $B$ 组成的序列作为输入时，它检查是否所有的 $A$ 都出现在所有的 $B$ 的左侧。该机器应将输入字符串保留在纸带上，并且当字符串是“按字母顺序”时停机，否则永远循环。
\end{prob}

\begin{prob}
\emph{字母顺序重排器：}设计一个字母表为 $\{\TMendtape,\TMblank, A, B\}$ 的Turing机，它以由有限个 $A$ 和 $B$ 组成的序列作为输入，并将它们重新排列，使得所有的 $A$ 都在所有的 $B$ 的左侧。（例如，序列 $BABAA$ 应变为 $AAABB$，序列 $ABBABB$ 应变为 $AABBBB$。）
\end{prob}

\end{document}